% ============================================================================
% Sección 3: Metodología
% ============================================================================
\section{Metodología}

\subsection{Enfoque de diseño}

El desarrollo de \texttt{search-library} siguió una metodología basada en principios de ingeniería de software, combinando diseño orientado a objetos con prácticas de desarrollo ágil. Se aplicaron los principios SOLID para garantizar extensibilidad y mantenibilidad del código \cite{martin2017clean}.

\subsubsection{Patrones de diseño aplicados}

\begin{itemize}
    \item \textbf{Strategy Pattern:} Permite intercambiar funciones heurísticas en tiempo de ejecución sin modificar el algoritmo base. Cada heurística implementa una interfaz común \texttt{Heuristic[T]}.
    
    \item \textbf{Adapter Pattern:} Los adaptadores \texttt{GraphSearchProblem} y \texttt{GridSearchProblem} transforman las estructuras de datos concretas (grafos, grids) en instancias compatibles con la interfaz abstracta \texttt{SearchProblem[T]}.
    
    \item \textbf{Template Method:} La clase abstracta \texttt{SearchProblem} define el esqueleto del problema de búsqueda, delegando la implementación específica a las subclases.
\end{itemize}

\subsection{Arquitectura del sistema}

La arquitectura de \texttt{search-library} se organiza en módulos cohesivos con responsabilidades bien definidas:

\begin{figure}[H]
\centering
\begin{tikzpicture}[
    module/.style={rectangle, draw=blue!60, fill=blue!5, thick, minimum width=2.8cm, minimum height=0.8cm, font=\small},
    arrow/.style={-{Stealth[length=3mm]}, thick},
    node distance=1.2cm and 1.5cm
]
    % Core module
    \node[module] (core) {\texttt{core}};
    \node[module, right=of core] (algorithms) {\texttt{algorithms}};
    \node[module, below=of core] (graph) {\texttt{graph}};
    \node[module, right=of graph] (grid) {\texttt{grid}};
    \node[module, below left=1cm and -0.5cm of graph] (heuristics) {\texttt{heuristics}};
    \node[module, below right=1cm and -0.5cm of grid] (exceptions) {\texttt{exceptions}};
    
    % Arrows
    \draw[arrow] (algorithms) -- (core);
    \draw[arrow] (graph) -- (core);
    \draw[arrow] (grid) -- (core);
    \draw[arrow] (grid) -- (heuristics);
    \draw[arrow] (graph) -- (heuristics);
    
\end{tikzpicture}
\caption{Arquitectura modular de \texttt{search-library}}
\label{fig:architecture}
\end{figure}

\subsection{Metodología experimental}

Para la validación del algoritmo se diseñó un protocolo experimental con los siguientes componentes:

\subsubsection{Escenarios de prueba}

\begin{enumerate}
    \item \textbf{Grafos ponderados:} Grafos dirigidos y no dirigidos con pesos variables, simulando redes de transporte con 10 a 100 nodos.
    
    \item \textbf{Grids bidimensionales:} Matrices de $10 \times 10$ hasta $50 \times 50$ con densidades de obstáculos del 20\% al 40\%, simulando entornos de navegación.
    
    \item \textbf{Configuraciones de heurística:} Evaluación comparativa entre heurística Manhattan (4 direcciones), Euclidiana (8 direcciones), y heurística nula ($h = 0$, equivalente a Dijkstra).
\end{enumerate}

\subsubsection{Métricas de evaluación}

Las métricas empleadas para la comparación son:

\begin{itemize}
    \item \textbf{Nodos explorados:} Número de estados expandidos durante la búsqueda.
    \item \textbf{Costo del camino:} Costo total de la solución encontrada.
    \item \textbf{Longitud del camino:} Número de pasos (aristas) en la solución.
    \item \textbf{Optimalidad:} Verificación de que la solución encontrada sea óptima.
\end{itemize}

\subsubsection{Herramientas y entorno}

\begin{itemize}
    \item \textbf{Lenguaje:} Python 3.11+
    \item \textbf{Framework de testing:} pytest con cobertura $\geq 90\%$
    \item \textbf{Análisis estático:} mypy (strict mode), ruff
    \item \textbf{Gestión de dependencias:} uv con hatchling como build system
    \item \textbf{Sin dependencias externas:} La librería no requiere paquetes de terceros
\end{itemize}

\subsection{Validación de correctitud}

La correctitud del algoritmo se valida mediante:

\begin{enumerate}
    \item \textbf{Pruebas unitarias:} Verificación de cada componente individual (nodos, grafos, grids, heurísticas).
    \item \textbf{Pruebas de integración:} Validación del flujo completo de búsqueda en escenarios conocidos.
    \item \textbf{Verificación de optimalidad:} Comparación de los resultados de A* contra soluciones calculadas por fuerza bruta en grafos pequeños.
    \item \textbf{Tipado estático:} Uso de \texttt{mypy} en modo estricto para garantizar corrección de tipos en tiempo de compilación.
\end{enumerate}
